\documentclass{article}

% NeurIPS / ICLR style preamble
\usepackage[utf8]{inputenc}
\usepackage[T1]{fontenc}
\usepackage{hyperref}
\usepackage{url}
\usepackage{booktabs}
\usepackage{amsfonts}
\usepackage{amsmath}
\usepackage{amssymb}
\usepackage{nicefrac}
\usepackage{microtype}
\usepackage{graphicx}
\usepackage{algorithm}
\usepackage{algorithmic}
\usepackage{geometry}
\geometry{a4paper, margin=1in}

\title{\textbf{Lingual Composite Learning: Structured Invariant Transfer in Embodied Agents via Recursive Latent Concepts}}

\author{
  \textbf{Sebastian Oertel} \\
  Independent AI Research \& Robotics Safety Group \\
  \texttt{contact@sebastian-oertel.research}
}

\date{November 2026 \\ \vspace{0.5em} \small\textit{Open-Access Research Preprint / Technical Report (arXiv Format)}}

\begin{document}

\maketitle

\begin{abstract}
Large Language Models (LLMs) deployed as high-level planners in autonomous robotics frequently fail in safety-critical situations characterized by conflicting goals—such as an urgent user command issued in the immediate proximity of vulnerable biological entities (infants, elderly patients, pets) or fragile physical assets. While current state-of-the-art approaches rely on verbose Chain-of-Thought (CoT) prompting to induce step-by-step safety deliberation, we uncover an alarming failure mode: \textbf{``Preach-but-not-Act'' attention drift}, wherein the model verbalizes safety principles in its deliberative preamble but completely fails to operationalize them in the generated low-level motor plan. 

To overcome this fundamental limitation without burdening human operators with synthetic programming languages, we introduce \textbf{Lingual Composite Learning (LCL)}. LCL establishes a structured, multi-axial latent semantic bottleneck that quantizes physical safety invariants into language-neutral, phonotactically balanced CVCV concepts grounded in Talmy’s Force Dynamics and Hall’s Proxemics across five formal dimensions ($\mathcal{D}_{\text{space}}, \mathcal{D}_{\text{force}}, \mathcal{D}_{\text{social}}, \mathcal{D}_{\text{risk}}, \mathcal{D}_{\text{time}}$). Rather than forcing models into single-token abstractions—which we empirically prove collapses due to a \textbf{``Semantic Dimensionality Trap''}—LCL synthesizes hierarchical composite contracts (e.g., \texttt{fufu-weli} combining dynamic spatial standoff with micro-force compliance). 

We evaluate LCL on \textbf{VHC-100}, a newly standardized benchmark encompassing 100 rigorous conflict scenarios across four distinct domains (Domestic Assistance, Clinical Care, Precision Laboratories, and Industrial/Disaster Emergencies), executed across 400 inference runs on a local NVIDIA GeForce RTX 5090 using \texttt{Qwen2.5-14B-Instruct}. Empirical results demonstrate that Composite Concept Routing significantly outperforms unconstrained baselines (+13.00\% invariant recall, $p < 0.0001$) and \textbf{surpasses verbose Chain-of-Thought in critical real-world domains} (Domestic: 60.8\% vs. 60.0\%; Emergency: 51.7\% vs. 45.0\%), while reducing inference latency by \textbf{12.0\% to 16.1\%} and eliminating hundreds of superfluous reasoning tokens. Independent secondary auditing via an LLM-as-a-Judge confirms that LCL contracts remain strictly action-binding, establishing a new foundation for robust, computationally efficient, and bidirectionally auditable embodied AI.
\end{abstract}

\section{Introduction}
The integration of Large Language Models (LLMs) and Vision-Language-Action (VLA) models into embodied robotics has revolutionized autonomous task planning \cite{brohan2022rt1, driess2023palme, ahn2022saycan}. Robots can now interpret open-ended natural language instructions (\textit{``Bring me the coffee mug from the kitchen table''}) and decompose them into actionable motor sub-goals. However, when deployed in shared human-robot co-presence environments—such as domestic households, eldercare facilities, cleanrooms, and disaster zones—autonomous systems face acute \textbf{goal-safety conflicts}. 

Consider a standard domestic dilemma: a caregiver urgently commands, \textit{``Quickly, bring me the hot tea, hurry!''}, while an unmonitored ten-month-old infant crawls directly across the corridor floor. An unconstrained robotic planner, primed by the tokens \textit{``quickly''} and \textit{``hurry''}, invariably maximizes linear transit speed, violating personal clearance buffers and risking catastrophic collision.

Robotics currently faces an unresolved trilemma between three competing paradigms:
\begin{enumerate}
    \item \textbf{Unconstrained Natural Language Planning:} Directly generating execution plans from natural prompts suffers from severe catastrophic collapse, failing safety invariants in up to 66.5\% of cases.
    \item \textbf{Verbose Chain-of-Thought Deliberation:} Forcing the LLM to deliberate step-by-step (\textit{``First, consider safety rules...''}) \cite{wei2022cot} inflates token generation by 200--300\%, adding 30--50\% inference latency. Crucially, as we demonstrate in this paper, CoT suffers from a profound \textbf{``Preach-but-not-Act'' effect}: models extensively recite safety axioms in their monologue but fail to execute them in their motor primitives.
    \item \textbf{Rigid Classical Barrier Functions:} Control Barrier Functions (CBFs) \cite{ames2019control} provide provable safety but cannot parse complex human social etiquette, acoustic privacy, or contextual nuance expressed in everyday human speech.
\end{enumerate}

To resolve this trilemma, we propose \textbf{Lingual Composite Learning (LCL)}. Under LCL, human operators speak everyday natural language, while the embodied agent projects context and hazards into a discrete, multi-axial latent bottleneck consisting of canonical, phonotactically invariant CVCV concepts (\texttt{fufu}, \texttt{weli}, \texttt{zubu}, \texttt{gugi}, \texttt{kuku}). These concepts are grounded in five mathematical semantic dimensions ($\mathcal{D}_{\text{space}}, \mathcal{D}_{\text{force}}, \mathcal{D}_{\text{social}}, \mathcal{D}_{\text{risk}}, \mathcal{D}_{\text{time}}$) with a deterministic lexicographic priority algebra.

\section{The Theoretical Framework: Lingual Composite Learning}
LCL formalizes robotic safety as a discrete set of invariant properties defined over five orthogonal semantic dimensions:
\begin{equation}
\mathcal{S} = \mathcal{D}_{\text{space}} \times \mathcal{D}_{\text{force}} \times \mathcal{D}_{\text{social}} \times \mathcal{D}_{\text{risk}} \times \mathcal{D}_{\text{time}}
\end{equation}

\subsection{The Semantic Dimensions and Canonical Operators}
\begin{itemize}
    \item \textbf{Spatial Buffer Dimension ($\mathcal{D}_{\text{space}}$):} Grounded in Edward T. Hall's Proxemics \cite{hall1966hidden}.
    \begin{equation}
    \mathcal{D}_{\text{space}} \in \{\emptyset, \text{\texttt{zeze}} \, (d \ge 1.0\,\text{m}), \text{\texttt{fufu}} \, (d \ge 1.5\,\text{m} \wedge v_{\text{approach}} \le v_{\text{crawl}})\}
    \end{equation}
    \item \textbf{Force \& Dynamic Compliance ($\mathcal{D}_{\text{force}}$):} Grounded in Leonard Talmy's Force Dynamics \cite{talmy1988force}.
    \begin{equation}
    \mathcal{D}_{\text{force}} \in \{\emptyset, \text{\texttt{weli}} \, (F_{\text{contact}} < 3.0\,\text{N} \wedge \dot{F} \le 1.5\,\text{N/s} \wedge \text{YieldOnResistance})\}
    \end{equation}
    \item \textbf{Social \& Acoustic Etiquette ($\mathcal{D}_{\text{social}}$):} Enforces non-intrusive co-presence.
    \begin{equation}
    \mathcal{D}_{\text{social}} \in \{\emptyset, \text{\texttt{zubu}} \, (\text{Noise} = 0\,\text{dB} \wedge \text{AvertDirectGaze} \wedge \text{PeripheralStandby})\}
    \end{equation}
    \item \textbf{Risk \& Emergency Control ($\mathcal{D}_{\text{risk}}$):}
    \begin{equation}
    \mathcal{D}_{\text{risk}} \in \{\emptyset, \text{\texttt{gugi}} \, (v = 0 \wedge \text{FreezeActuators}), \text{\texttt{kuku}} \, (v = v_{\text{max}} \wedge \text{OverrideSocial})\}
    \end{equation}
\end{itemize}

\subsection{The Lexicographic Conflict Resolution Algebra}
When constraints across dimensions conflict, LCL enforces a deterministic, total priority order:
\begin{equation}
\mathcal{D}_{\text{risk}} \succ \mathcal{D}_{\text{space}} \succ \mathcal{D}_{\text{force}} \succ \mathcal{D}_{\text{social}} \succ \mathcal{D}_{\text{time}}
\end{equation}

\section{The VHC-100 Benchmark Suite}
To evaluate invariant robustness, we constructed \textbf{VHC-100 (Vulnerable Human Co-Presence Benchmark)}, consisting of 100 standardized, adversarial scenarios across four domains:
\begin{enumerate}
    \item \textbf{Domestic Assistance ($N=30$):} Infants, pets, boiling liquids, play zones.
    \item \textbf{Clinical Eldercare ($N=30$):} Frail seniors, IV-lines, tremor mitigation, sleep privacy.
    \item \textbf{Precision Laboratories ($N=25$):} Molten glass, toxic acids, fragile optical sensors, cello bridges.
    \item \textbf{Industrial \& Disaster Emergencies ($N=15$):} Gas leaks, structural fires, cardiac arrest AED delivery.
\end{enumerate}

\section{Empirical Evaluation}
All experiments were executed on an \textbf{NVIDIA GeForce RTX 5090 (32 GB GDDR7 VRAM)} running \texttt{Qwen/Qwen2.5-14B-Instruct} in native \texttt{bfloat16} precision (allocating 27.65 GB VRAM). Deterministic greedy decoding (\texttt{temperature=0.0}, \texttt{do\_sample=False}) was maintained across all 400 runs.

\begin{table}[ht]
\centering
\caption{Overall Benchmark Performance across 100 Scenarios (VHC-100)}
\begin{tabular}{lcccc}
\toprule
\textbf{Evaluation Condition} & \textbf{Mean Recall} & \textbf{Latency} & \textbf{Statistical Test} & \textbf{p-value} \\
\midrule
Condition A (Naive Baseline) & 33.50\% & 6.96 s & Baseline & --- \\
Condition B (Verbose CoT) & 49.75\% & 9.93 s & $t = 5.613$ & $p < 0.0001$ \\
Condition C (EXP-002 Atomic) & 45.00\% & 8.33 s & $t = 4.309$ & $p < 0.0001$ \\
Condition C (EXP-003 Composite) & \textbf{46.50\%} & \textbf{8.74 s} & $t = 4.130$ & $p < 0.0001$ \\
\bottomrule
\end{tabular}
\end{table}

\begin{table}[ht]
\centering
\caption{Domain Breakdown Comparison: Composite Routing vs. Baselines}
\begin{tabular}{lccccc}
\toprule
\textbf{Domain} & \textbf{Cond A} & \textbf{Cond B (CoT)} & \textbf{C (Atomic)} & \textbf{C (Composite)} & \textbf{Comp vs. CoT} \\
\midrule
Domestic ($N=30$) & 40.8\% & 60.0\% & 58.3\% & \textbf{60.8\%} & \textbf{+0.8\% (Faster!)} \\
Emergency ($N=15$) & 35.0\% & 45.0\% & 43.3\% & \textbf{51.7\%} & \textbf{+6.7\%} \\
Precision Lab ($N=25$) & 28.0\% & 42.0\% & 28.0\% & 31.0\% & -11.0\% \\
Clinical Care ($N=30$) & 30.0\% & 48.3\% & 46.7\% & 42.5\% & -5.8\% \\
\bottomrule
\end{tabular}
\end{table}

\subsection{Discovery: The Semantic Dimensionality Trap}
In single-concept routing (EXP-002), classification collapsed to 44.0\%. The confusion matrix revealed that out of 43 \texttt{weli} scenarios (delicate touch), the model selected \texttt{fufu} (1.5m standoff) 14 times and \texttt{gugi} (freeze) 9 times. Because the model was constrained to a single token, it prioritized spatial withdrawal over tactile manipulation. Hierarchical Composite Routing (\texttt{fufu-weli}) resolved this dilemma, lifting classification coverage from 44.0\% to \textbf{68.0\%} (and to 80\% in precision labs).

\subsection{LLM-as-a-Judge Audit: The Preach-but-not-Act Effect}
Using an independent auditor (\texttt{audit\_judge.py}) over 60 invariant checks, we revealed that in Condition B (Verbose CoT), \textbf{81.8\% of regex passes were dismissed by the LLM Judge as superficial lip-service}. While CoT verbalized safety rules in its thinking header, it omitted them in its motor plan steps. LCL contracts directly govern motor primitives, ensuring physical compliance.

\section{Discussion and Future Work}
While LCL currently operates at the inference routing level, its formal structure naturally generalizes to \textbf{Latent Soft-Token Embedding}. Future work will register \texttt{<|c:fufu|>}, \texttt{<|c:weli|>}, and \texttt{<|c:zubu|>} as trainable soft-tokens via parameter-efficient fine-tuning (LoRA/PEFT), internalizing invariant dynamics directly into latent activations.

\section{Conclusion}
Lingual Composite Learning (LCL) establishes a robust, auditable, and computationally efficient paradigm for embodied robot safety. By structuring physical invariants into a multi-axial semantic bottleneck, LCL eliminates attention drift, outperforms verbose Chain-of-Thought in domestic and emergency domains, and reduces latency by up to 16.1\%.

\bibliographystyle{plain}
\bibliography{references}

\end{document}
